\SourceLeaf{156}%
Hinc fit
\[
\varphi(2b,a-b)=2\varphi(b,a)-\varphi(2b,a)+\frac14b(a-3-b)
\]
et proin
\[
\varphi(a-b,a+b)=\varphi(2b,a)-2\varphi(b,a)+\frac18(aa-bb+2b-1)
\]
Substituendo hunc valorem in formula pro $g$ supra tradita, insuperque $\varphi(a,b)+\varphi(b,a)=\frac14b(a-1)$, obtinemus
\[
g=2\varphi(2b,a)-2\varphi(b,a)+\frac18(aa-2ab+bb+4b-1)
\]

\Article{74}
Per ratiocinia prorsus similia absolvitur casus is, ubi manentibus $a$, $b$ positivis $a-b$ est negativus, sive $b-a$ positivus. Aequationes $p(a-2x)=2b\eta+a\theta$, $p(b-a+2x)=2b\zeta+(b-a)\theta$ docent, $\frac12a-x$ atque $x+\frac12(b-a)$ positivos, et proin $x$ alicui numerorum $-\frac12(b-a-1),-\frac12(b-a-3),-\frac12(b-a-5)\dotsfour+\frac12(a-1)$ aequalem esse debere. Porro ex aequatione $px+(b-a)\eta=a\zeta$ sequitur, pro valoribus negativis ipsius $x$ conditionem, ex qua $\eta$ debet esse positivus, iam contineri sub conditione, ex qua $\zeta$ debet esse positivus, contrarium vero evenire, quoties ipsi $x$ valor positivus tribuatur. Hinc valores ipsius $y$ pro valore determinato negativo ipsius $x$ inter $\frac{(a+b)x}{a-b}$ et $\frac{p-2ax}{2b}$, contra pro valore positivo ipsius $x$ inter $\frac{bx}{a}$ et $\frac{p-2ax}{2b}$ contenti esse debent: manifesto pro $x=0$ hi limites sunt $0$ et $\frac{p-2ax}{2b}$, valore $y=0$ ipso excluso. Hinc colligitur
\[
g=-\Sigma\left[\frac{(a+b)x}{a-b}\right]+\Sigma\left[\frac12b+\frac{aa-2ax}{2b}\right]-\Sigma\left[\frac{bx}{a}\right]
\]
ubi in termino primo summatio extendenda est per omnes valores negativos ipsius. $x$ inde a $-1$ usque ad $-\frac12(b-a-1)$; in secunda per omnes valores ipsius $x$ inde a $-\frac12(b-a-1)$ usque ad $\frac12(a-1)$; in tertia per omnes valores positivos ipsius $x$ inde a $+1$ usque ad $\frac12(a-1)$: hoc pacto e summatione prima prodit $-\varphi(b-a,b+a)$, e secunda perinde ut in art. praec. $\frac14bb+\varphi(2b,a)-\varphi(b,a)$, denique e tertia $-\varphi(a,b)$, sive habetur
\[
g=-\varphi(b-a,b+a)+\varphi(2b,a)-\varphi(b,a)-\varphi(a,b)+\frac14bb
\]
Iam simili modo ut in art. praec. evolvitur
\SourceLeaf{157}%
\[
\begin{aligned}
\varphi(b-a,b+a)&=\varphi(b-a,2b)-\frac18((b-a)^2-1)\\
&=\frac18(3bb-2ab-aa-4b+1)-\varphi(2b,b-a)
\end{aligned}
\]
nec non
\[
\varphi(2b,b-a)=\varphi(2b,a)-2\varphi(b,a)+\frac14b(b-1-a)
\]
adeoque
\[
\varphi(b-a,b+a)=2\varphi(b,a)-\varphi(2b,a)+\frac18(bb-aa-2b+1)
\]
tandemque
\[
g=2\varphi(2b,a)-2\varphi(b,a)+\frac18(aa-2ab+bb+4b-1)
\]
Evictum est itaque, eandem formulam pro $g$ valere, sive sit $a-b$ positivus sive negativus, dummodo $a$, $b$ sint positivi.

\Article{75}
Ut reductionem ulteriorem assequamur, statuemus
\[
\begin{aligned}
L&=\left[\frac{a}{2b}\right]+\left[\frac{2a}{2b}\right]+\left[\frac{3a}{2b}\right]+\text{etc.}+\left[\frac{\frac12ba}{2b}\right]\\
M&=\left[\frac{(\frac12b+1)a}{2b}\right]+\left[\frac{(\frac12b+2)a}{2b}\right]+\left[\frac{(\frac12b+3)a}{2b}\right]+\text{etc.}+\left[\frac{ba}{2b}\right]\\
N&=\left[\frac{a+b}{2b}\right]+\left[\frac{2a+b}{2b}\right]+\left[\frac{3a+b}{2b}\right]+\text{etc.}+\left[\frac{\frac12ba+b}{2b}\right]
\end{aligned}
\]
Quum facile perspiciatur, haberi generaliter $[u]+[u+\frac12]=[2u]$, quamcunque quantitatem realem denotet $u$, fit $L+N=\varphi(b,a)$, et quum manifesto sit $L+M=\varphi(2b,a)$, erit
\[
\varphi(2b,a)-\varphi(b,a)=M-N
\]

Porro autem obvium est, aggregatum termini primi seriei $N$ cum penultimo termino seriei $M$, puta $\left[\frac{a+b}{2b}\right]+\left[\frac{(b-1)a}{2b}\right]$ fieri $=\frac12(a-1)$, atque eandem summam effici e termino secundo seriei $N$ cum antepenultimo seriei $M$, et sic porro: quare quum etiam terminus ultimus seriei $M$ fiat $=\frac12(a-1)$, ultimus vero terminus seriei $N$ sit $=\left[\frac{a+2}{4}\right]=\frac14(a\mp1)$, valente signo superiori vel inferiori, prout $a$ est formae $4n+1$ vel $4n-1$: erit
\[
M+N=\frac14(a-1)b+\frac14(a\mp1)
\]
et proin
\[
\varphi(2b,a)-\varphi(b,a)=\frac14(a-1)b+\frac14(a\mp1)-2N
\]

\SourceLeaf{158}%
Formula itaque pro $g$ in artt. 73 et 74 inventa, transit in sequentem
\[
g=\frac18((a+b)^2-1)+2n-4N
\]
statuendo $a\mp1=4n$, ubi $n$ erit integer. Sed quum hinc habeatur $1=16nn-8an+aa$, formula haec etiam sequenti modo exhiberi potest:
\[
g=\frac18(-aa+2ab+bb+1)+4(\frac12(a+1)n-nn-N)
\]
Quapropter quum $g$ sit character numeri $-1-i$ pro modulo $a+bi$, hic character fit $\equiv\frac18(-aa+2ab+bb+1)\pmod4$, quod est ipsum theorema supra (art. 64) per inductionem erutum, sponteque inde demanant theoremata circa characteres numerorum $1+i$, $1-i$, $-1+i$. Quamobrem haec quatuor theoremata, pro casu eo, ubi $a$ et $b$ sunt positivi, iam rigorose sunt demonstrata.

\Article{76}
Si manente $a$ positivo $b$ est negativus, statuatur $b=-b'$, ut fiat $b'$ positivus. Quum iam evictum sit, ita pro modulo $a+b'i$ characterem numeri $-1-i$ esse $\equiv\frac18(-aa+2ab'+b'b'+1)\pmod4$, character numeri $-1+i$ pro modulo $a-b'i$ per theorema in art. 62 prolatum erit $\equiv\frac18(aa-2ab'-b'b'-1)$, i.e. character numeri $-1+i$ pro modulo $a+bi$ fit $\equiv\frac18(aa+2ab-bb-1)$: hoc vero est ipsum theorema in art. 64 allatum, unde tria reliqua circa characteres numerorum $1+i$, $1-i$, $-1-i$ sponte demanant. Quapropter ista theoremata etiam pro casu, ubi $b$ negativus est, demonstrata sunt, scilicet pro omnibus casibus, ubi $a$ est positivus.

Denique si $a$ est negativus, statuatur $a=-a'$, $b=-b'$. Quum itaque per iam demonstrata character numeri $1+i$ respectu moduli $a'+b'i$ sit $\equiv\frac18(-a'a'+2a'b'-3b'b'+1)\pmod4$, nihilque intersit, utrum numerum $a'+b'i$ an oppositum $-a'-b'i$ moduli loco habeamus; manifesto character numeri $1+i$ respectu moduli $a+bi$ est $\equiv\frac18(-aa+2ab-3bb+1)$, et similia valent circa characteres numerorum $1-i$, $-1+i$, $-1-i$. 

Ex his itaque colligitur, demonstrationem theorematum circa characteres numerorum $1+i$, $1-i$, $-1+i$, $-1-i$ (artt. 63. 64) nulli amplius limitationi obnoxiam esse.
\BookRule{.28\linewidth}

\clearpage\SourceLeaf{159}%
\thispagestyle{empty}
\pdfbookmark[0]{Anzeigen eigner Schriften}{notices-title}
\begin{center}
\vspace*{1in}{\Large ANZEIGEN\par}
\vspace{.4in}{\large EIGNER\par}
\vspace{.4in}{\Large SCHRIFTEN.\par}
\end{center}

